\begin{filecontents*}[overwrite]{\jobname.ist}
headings_flag 1
numhead_positive "数字"
symhead_positive "符号"
heading_prefix "{\\bfseries\\hfil "
heading_suffix "\\hfil}\\nopagebreak\\par\\vskip 5pt "
\end{filecontents*}

\setCJKmainfont[
  ItalicFont     = KaiTi,
  BoldFont       = SimHei,
  BoldItalicFont = FangSong,
]{SimSun}

\usepackage{amsmath,amssymb,amsthm}
\usepackage{mathtools}
\usepackage{tikz}              %
\usepackage{etoolbox}          %
\usepackage[a4paper,margin=2.8cm]{geometry}
\usepackage{enumitem}
\usepackage{longtable}  %
\usepackage{caption}    %
\captionsetup{labelsep=quad}
\setlist{itemsep=2pt,topsep=4pt}
\emergencystretch=1.5em %
\vfuzz=3pt              %
\hbadness=2500          %
\xeCJKsetup{PunctStyle=quanjiao}%

\newif\ifzhnames
\zhnamestrue   %

\newcommand{\trInfinityCategory}{∞-范畴}%
\let\trInfinityCategories\trInfinityCategory%
\newcommand{\trQuasiCategory}{拟范畴}%
\newcommand{\trCategory}{范畴}%
\let\trCategories\trCategory%
\newcommand{\trSimplicialSet}{单纯集}%
\let\trSimplicialSets\trSimplicialSet%
\newcommand{\trSimplicialCategory}{单纯范畴}%
\let\trSimplicialCategories\trSimplicialCategory%
\newcommand{\trNerve}{神经}%
\newcommand{\trHorn}{角}%
\newcommand{\trSimplex}{单形}%
\newcommand{\trKanComplex}{Kan 复形}%
\newcommand{\trKanFibration}{Kan 纤维化}%
\newcommand{\trWeakEquivalence}{弱等价}%
\newcommand{\trJoyalEquivalence}{若亚尔等价}%
\newcommand{\trMappingSpace}{映射空间}%
\newcommand{\trFunctor}{函子}%
\newcommand{\trNaturalTransformation}{自然变换}%
\newcommand{\trEquivalence}{等价}%
\newcommand{\trLimit}{极限}%
\newcommand{\trColimit}{余极限}%
\newcommand{\trTerminalObject}{终对象}%
\newcommand{\trInitialObject}{始对象}%
\newcommand{\trAdjunction}{伴随}%
\newcommand{\trLeftAdjoint}{左伴随}%
\newcommand{\trRightAdjoint}{右伴随}%
\newcommand{\trAdjointFunctorTheorem}{伴随函子定理}%
\newcommand{\trLocalization}{局部化}%
\newcommand{\trFibration}{纤维化}%
\newcommand{\trInnerFibration}{内纤维化}%
\newcommand{\trInnerAnodyne}{内零调}%
\newcommand{\trAnodyneMap}{零调映射}%
\newcommand{\trCartesianFibration}{笛卡尔纤维化}%
\newcommand{\trCocartesianFibration}{余笛卡尔纤维化}%
\newcommand{\trMarkedSimplicialSet}{带标记单纯集}%
\newcommand{\trMarkedAnodyne}{带标记零调}%
\newcommand{\trStraightening}{拉直}%
\newcommand{\trUnstraightening}{反拉直}%
\newcommand{\trTwistedArrowCategory}{扭箭范畴}%
\newcommand{\trJoin}{联}%
\newcommand{\trFatJoin}{胖联}%
\newcommand{\trFatSlice}{胖切片}%
\newcommand{\trHomotopyEquivalence}{同伦等价}%
\newcommand{\trSliceCategory}{切片范畴}%
\newcommand{\trCofinal}{共尾的}%
\newcommand{\trCoinitial}{共始的}%
\newcommand{\trFullyFaithful}{全忠实的}%
\newcommand{\trEssentiallySurjective}{本质满的}%
\newcommand{\trHomotopy}{同伦}%
\newcommand{\trFiltered}{滤过的}%
\newcommand{\trCategoryTheory}{范畴论}%
\newcommand{\trHigherCategoryTheory}{高阶范畴论}%
\newcommand{\trHomotopyTheory}{同伦论}%
\newcommand{\trModelCategory}{模型范畴}%
\newcommand{\trAlgebraicTopology}{代数拓扑}%
\newcommand{\trTopologicalSpace}{拓扑空间}%
\newcommand{\trEnrichment}{充实}%
\newcommand{\trCoherentNerve}{凝聚神经}%
\newcommand{\trObject}{对象}%
\newcommand{\trMorphism}{态射}%
\newcommand{\trMonomorphism}{单态射}%
\newcommand{\trPresheaf}{预层}%
\newcommand{\trGroupoid}{群胚}%
\newcommand{\trPushout}{推出}%
\newcommand{\trPullback}{拉回}%
\newcommand{\trKanExtension}{Kan 扩张}%
\newcommand{\trSpine}{脊}%
\newcommand{\trSkeleton}{骨架}%
\newcommand{\trCoskeleton}{余骨架}%
\newcommand{\trGeometricRealization}{几何实现}%
\newcommand{\trSingularSimplicialSet}{奇异单纯集}%
\newcommand{\trCosimplicial}{余单纯}%
\newcommand{\trFaceMap}{面映射}%
\newcommand{\trDegeneracyMap}{退化映射}%
\newcommand{\trBicomplete}{双完备}%
\newcommand{\trCocone}{余锥}%
\newcommand{\trLiftingProblem}{提升问题}%
\newcommand{\trComposer}{合成器}%
\let\trWeakComposer\trComposer%
\newcommand{\trMonoidalStructure}{幺半结构}%
\newcommand{\trMonoidalProduct}{幺半积}%
\newcommand{\trMonoidalAdjunction}{幺半伴随}%
\newcommand{\trLaxMonoidalFunctor}{松幺半函子}%
\newcommand{\trOplaxMonoidalFunctor}{反松幺半函子}%
\newcommand{\trUnderlyingCategory}{底范畴}%
\newcommand{\trSimplicialNerve}{单纯神经}%
\newcommand{\trInnerFace}{内面}%
\newcommand{\trPuncturedCube}{穿孔立方体}%
\newcommand{\trTwoFunctor}{2-函子}%
\newcommand{\trTwoMorphism}{2-态射}%
\newcommand{\trCWComplex}{CW 复形}%
\newcommand{\trWeaklyFullyFaithful}{弱全忠实的}%
\newcommand{\trWeaklyEssentiallySurjective}{弱本质满的}%
\newcommand{\trLeftFibration}{左纤维化}%
\newcommand{\trRightFibration}{右纤维化}%
\newcommand{\trTrivialFibration}{平凡纤维化}%
\newcommand{\trAnodyne}{零调}%
\newcommand{\trLeftAnodyne}{左零调}%
\newcommand{\trRightAnodyne}{右零调}%
\newcommand{\trSaturated}{饱和的}%
\newcommand{\trSaturatedSet}{饱和集}%
\newcommand{\trSaturation}{饱和闭包}%
\newcommand{\trRetract}{收缩}%
\newcommand{\trRightLiftingProperty}{右提升性质}%
\newcommand{\trLeftLiftingProperty}{左提升性质}%
\newcommand{\trSmallObjectArgument}{小对象论证}%
\newcommand{\trBoundary}{边界}%
\newcommand{\trHornInclusion}{角包含}%
\newcommand{\trBoundaryInclusion}{边界包含}%
\newcommand{\trSpineInclusion}{脊包含}%
\newcommand{\trCompact}{紧的}%
\newcommand{\trSectionMap}{截面}%
\newcommand{\trPushoutProduct}{推出积}%
\newcommand{\trJoyalFibration}{若亚尔纤维化}%
\newcommand{\trModelStructure}{模型结构}%
\newcommand{\trClassifier}{分类器}%
\newcommand{\trCorepresented}{余表示的}%
\newcommand{\trSimplicialMap}{单纯映射}%
\newcommand{\trInfinityGroupoid}{∞-群胚}%
\newcommand{\trCut}{切割}%
\newcommand{\trSlice}{切片}%
\newcommand{\trCone}{锥}%
\newcommand{\trLinearlyOrderedSet}{线性有序集}%
\newcommand{\trForgetfulFunctor}{遗忘函子}%
\newcommand{\trContravariantFunctor}{反变函子}%
\newcommand{\trRepresentable}{可表示的}%
\newcommand{\trComplete}{完备的}%
\newcommand{\trCocomplete}{余完备的}%
\newcommand{\trCoequalizer}{余等化子}%
\newcommand{\trUnit}{单位}%
\newcommand{\trCounit}{余单位}%
\newcommand{\trConnected}{连通}%
\newcommand{\trSaturatedClass}{饱和类}%
\newcommand{\trConservative}{保守的}%
\newcommand{\trOppositeFunctor}{相反函子}%
\newcommand{\trMonoidalFunctor}{幺半函子}%
\newcommand{\trCategoricalEquivalence}{范畴等价}%
\newcommand{\trNaturalEquivalence}{自然等价}%
\newcommand{\trIsofibration}{等价纤维化}%
\newcommand{\trPointwise}{逐点}%
\newcommand{\trInverseLimit}{逆极限}%
\newcommand{\trFibre}{纤维}%
\newcommand{\trOrdinaryCategory}{通常范畴}%
\newcommand{\trThreeForTwoProperty}{三取二性质}%
\newcommand{\trLoop}{闭路}%
\newcommand{\trTwoCell}{2-胞腔}%
\newcommand{\trInvertible}{可逆的}%
\newcommand{\trHomotopyCategory}{同伦范畴}%
\newcommand{\trFibrantObject}{纤维性对象}%
\newcommand{\trFibreSequence}{纤维列}%
\newcommand{\trHeuts}{\ifzhnames 赫伊茨\else Heuts\fi}%
\newcommand{\trWhitehead}{\ifzhnames 怀特海\else Whitehead\fi}%
\newcommand{\trReedy}{\ifzhnames 里迪\else Reedy\fi}%
\newcommand{\trEssentialImage}{本质像}%
\newcommand{\trUniversalProperty}{泛性质}%
\newcommand{\trPathFibration}{路纤维化}%
\newcommand{\trFiltration}{滤过}%
\newcommand{\trContractible}{可缩的}%
\newcommand{\trOuterHorn}{外角}%
\newcommand{\trGrothendieck}{\ifzhnames 格罗滕迪克\else Grothendieck\fi}%
\newcommand{\trGrothendieckConstruction}{格罗滕迪克构造}%
\newcommand{\trPseudoFunctor}{伪函子}%
\newcommand{\trTwoCategory}{2-范畴}%
\newcommand{\trMarkedLeftAnodyne}{带标记左零调的}%
\newcommand{\trLocallyCocartesian}{局部余笛卡尔的}%
\newcommand{\trMarkedEquivalence}{带标记等价}%
\newcommand{\trMarkedEdge}{带标记边}%
\newcommand{\trCartesianClosed}{笛卡尔闭的}%
\newcommand{\trSimplicialModelCategory}{单纯模型范畴}%
\newcommand{\trCofibration}{上纤维化}%
\newcommand{\trCofibrant}{上纤维的}%
\newcommand{\trInitialVertexMap}{始顶点映射}%
\newcommand{\trUniversalCocartesianFibration}{泛余笛卡尔纤维化}%
\newcommand{\trHomotopyCartesian}{同伦笛卡尔的}%
\newcommand{\trStevenson}{\ifzhnames 史蒂文森\else Stevenson\fi}%
\newcommand{\trConstantFunctor}{常值函子}%
\newcommand{\trOppositeCategory}{相反范畴}%
\newcommand{\trExtensionProblem}{扩张问题}%
\newcommand{\trBivariant}{双变的}%
\newcommand{\trHebestreit}{\ifzhnames 赫贝斯特赖特\else Hebestreit\fi}%
\newcommand{\trCoproduct}{余积}%
\newcommand{\trProduct}{积}%
\newcommand{\trDiscreteCategory}{离散范畴}%
\newcommand{\trRepresentingObject}{表示对象}%
\newcommand{\trWeaklyContractible}{弱可缩的}%
\newcommand{\trKrause}{\ifzhnames 克劳泽\else Krause\fi}%
\newcommand{\trNguyen}{\ifzhnames 阮黄金\else Hoang Kim Nguyen\fi}%
\newcommand{\trCofinality}{共尾性}%
\newcommand{\trSmooth}{光滑的}%
\newcommand{\trProper}{正常的}%
\newcommand{\trUniversallySmooth}{泛光滑的}%
\newcommand{\trUniversallyProper}{泛正常的}%
\newcommand{\trBicartesianFibration}{双笛卡尔纤维化}%
\newcommand{\trBinaturalTransformation}{双自然变换}%
\newcommand{\trBousfield}{Bousfield}%
\newcommand{\trDiagonalFunctor}{对角函子}%
\newcommand{\trEssentiallySmall}{本质小的}%
\newcommand{\trLocallySmall}{局部小的}%
\newcommand{\trWeaklyInitial}{弱始的}%
\newcommand{\trNonDegenerate}{非退化的}%
\newcommand{\trSimplicialIdentities}{单纯恒等式}%
\newcommand{\trRetraction}{收缩映射}%
\newcommand{\trPointed}{带点的}%
\newcommand{\trInitialCone}{始锥}%
\newcommand{\trPresentable}{可表现的}%
\newcommand{\trAccessible}{可及的}%
\newcommand{\trWeaklyTerminal}{弱终的}%

\newcommand{\trhTheorem}{定理}%
\newcommand{\trhProposition}{命题}%
\newcommand{\trhLemma}{引理}%
\newcommand{\trhCorollary}{推论}%
\newcommand{\trhFact}{事实}%
\newcommand{\trhDefinition}{定义}%
\newcommand{\trhExample}{例}%
\newcommand{\trhNotation}{记号}%
\newcommand{\trhConvention}{约定}%
\newcommand{\trhConstruction}{构造}%
\newcommand{\trhRemark}{注}%
\newcommand{\trhObservation}{观察}%
\newcommand{\trhAxiom}{公理}%
\newcommand{\trhWarning}{警告}%
\newcommand{\trhExercise}{练习}%
\newcommand{\trhProof}{证明}%

\newcommand{\trLand}{\ifzhnames 兰德\else Land\fi}%
\newcommand{\trJoyal}{\ifzhnames 若亚尔\else Joyal\fi}%
\newcommand{\trLurie}{\ifzhnames 卢瑞\else Lurie\fi}%
\newcommand{\trQuillen}{\ifzhnames 奎伦\else Quillen\fi}%
\newcommand{\trDwyer}{\ifzhnames 德怀尔\else Dwyer\fi}%
\newcommand{\trKan}{Kan}%
\newcommand{\trCisinski}{\ifzhnames 西桑斯基\else Cisinski\fi}%
\newcommand{\trYoneda}{\ifzhnames 米田\else Yoneda\fi}%
\newcommand{\trRezk}{\ifzhnames 雷兹\else Rezk\fi}%
\newcommand{\trHaugseng}{\ifzhnames 豪格森\else Haugseng\fi}%
\newcommand{\trGepner}{\ifzhnames 盖普纳\else Gepner\fi}%
\newcommand{\trNikolaus}{\ifzhnames 尼古拉斯\else Nikolaus\fi}%
\newcommand{\trGalatius}{\ifzhnames 加拉提乌斯\else Galatius\fi}%

\newcommand{\trSource}{源}%
\newcommand{\trTarget}{靶}%
\newcommand{\trGoodwillie}{\ifzhnames 古德威利\else Goodwillie\fi}%
\newcommand{\trGoerss}{\ifzhnames 格斯\else Goerss\fi}%
\newcommand{\trHopkins}{\ifzhnames 霍普金斯\else Hopkins\fi}%
\newcommand{\trMiller}{\ifzhnames 米勒\else Miller\fi}%
\newcommand{\trSegal}{\ifzhnames 西格尔\else Segal\fi}%
\newcommand{\trTate}{\ifzhnames 泰特\else Tate\fi}%
\newcommand{\trPoincare}{\ifzhnames 庞加莱\else Poincaré\fi}%
\newcommand{\trAtiyah}{\ifzhnames 阿蒂亚\else Atiyah\fi}%
\newcommand{\trSpectrum}{谱}%
\newcommand{\trSheaf}{层}%
\newcommand{\trDescent}{下降}%
\newcommand{\trGroupCompletion}{群完备化}%
\newcommand{\trScheme}{概形}%
\newcommand{\trDerivedCategory}{导出范畴}%
\newcommand{\trSmashProduct}{打积}%
\newcommand{\trMotivic}{母题的}%
\newcommand{\trLocalSystem}{局部系统}%
\newcommand{\trSurgeryTheory}{割补理论}%
\newcommand{\trFullSubcategory}{满子范畴}%
\newcommand{\trClassifyingMap}{分类映射}%
\newcommand{\trHomotopyHypothesis}{同伦假设}%
\newcommand{\trFreeStandingIsomorphism}{自由驻立同构}%
\newcommand{\trLocallyConstant}{局部常值的}%
\newcommand{\trHermitian}{埃尔米特}%
\newcommand{\trComposable}{可合成的}%


\newcounter{icitem}[section]
\renewcommand{\theicitem}{\thesection.\arabic{icitem}}
\newcounter{icex}   %
\renewcommand{\theicex}{\arabic{icex}}

\newtheoremstyle{icplain}{\topsep}{\topsep}{\itshape}{}
  {\bfseries}{}{.5em}{#1~#2\ifstrempty{#3}{}{（{\kaishu #3}）}}
\newtheoremstyle{icdefinition}{\topsep}{\topsep}{}{}
  {\bfseries}{}{.5em}{#1~#2\ifstrempty{#3}{}{（{\kaishu #3}）}}
\newtheoremstyle{icremark}{\topsep}{\topsep}{\upshape}{}
  {\itshape}{}{.5em}{#1~#2\ifstrempty{#3}{}{（{\kaishu #3}）}}

\theoremstyle{icplain}
\newtheorem{theorem}[icitem]{\trhTheorem}
\newtheorem{proposition}[icitem]{\trhProposition}
\newtheorem{lemma}[icitem]{\trhLemma}
\newtheorem{corollary}[icitem]{\trhCorollary}
\newtheorem{fact}[icitem]{\trhFact}
\theoremstyle{icdefinition}
\newtheorem{definition}[icitem]{\trhDefinition}
\newtheorem{example}[icitem]{\trhExample}
\newtheorem{notation}[icitem]{\trhNotation}
\newtheorem{convention}[icitem]{\trhConvention}
\newtheorem{construction}[icitem]{\trhConstruction}
\theoremstyle{icremark}
\newtheorem{remark}[icitem]{\trhRemark}
\newtheorem{observation}[icitem]{\trhObservation}
\newtheorem{warning}[icitem]{\trhWarning}
\newtheorem*{warning*}{\trhWarning}

\theoremstyle{icdefinition}
\newtheorem*{example*}{\trhExample}
\theoremstyle{icplain}
\newtheorem*{axiom*}{\trhAxiom}

\newtheorem{exercise}[icex]{\trhExercise}

\newcommand{\icsync}[1]{\setcounter{icitem}{\numexpr(#1)-1\relax}}
\newcommand{\icsyncex}[1]{\setcounter{icex}{\numexpr(#1)-1\relax}}

\newenvironment{iceq}{\begin{equation*}\refstepcounter{icitem}\tag{\theicitem}}{\end{equation*}}

\renewcommand{\proofname}{\trhProof}
\renewcommand{\qedsymbol}{$\blacksquare$}

\counterwithin{figure}{chapter}
\counterwithin{table}{chapter}

\ctexset{
  chapter = {
    name       = {第,章},
    number     = {\arabic{chapter}},
    format     = {\huge\bfseries\raggedright},
    nameformat = {\large\bfseries},
    numberformat= {\large\bfseries},
    aftername  = {\par\vspace{14pt}},
    beforeskip = 24pt,
    afterskip  = 40pt,
  },
}

\AtBeginEnvironment{theorem}{\crefalias{icitem}{theorem}}
\AtBeginEnvironment{proposition}{\crefalias{icitem}{proposition}}
\AtBeginEnvironment{lemma}{\crefalias{icitem}{lemma}}
\AtBeginEnvironment{corollary}{\crefalias{icitem}{corollary}}
\AtBeginEnvironment{fact}{\crefalias{icitem}{fact}}
\AtBeginEnvironment{definition}{\crefalias{icitem}{definition}}
\AtBeginEnvironment{example}{\crefalias{icitem}{example}}
\AtBeginEnvironment{notation}{\crefalias{icitem}{notation}}
\AtBeginEnvironment{convention}{\crefalias{icitem}{convention}}
\AtBeginEnvironment{construction}{\crefalias{icitem}{construction}}
\AtBeginEnvironment{remark}{\crefalias{icitem}{remark}}
\AtBeginEnvironment{observation}{\crefalias{icitem}{observation}}
\AtBeginEnvironment{warning}{\crefalias{icitem}{warning}}
\AtBeginEnvironment{exercise}{\crefalias{icex}{exercise}}

\makeatletter
\patchcmd{\@oddhead}{\slshape}{}{}{}
\patchcmd{\@evenhead}{\slshape}{}{}{}
\makeatother

\newcommand{\R}{\mathbb{R}}
\newcommand{\N}{\mathbb{N}}
\newcommand{\Z}{\mathbb{Z}}
\newcommand{\C}{\mathcal{C}}   %
\newcommand{\D}{\mathcal{D}}
\newcommand{\E}{\mathcal{E}}
\newcommand{\Q}{\mathbb{Q}}
\newcommand{\oo}{\ensuremath{\infty}}  %

\usepackage{xcolor}
\usepackage{hyperref}
\hypersetup{
  colorlinks=true,
  linkcolor=blue!55!black,
  citecolor=teal,
  urlcolor=violet,
  bookmarksnumbered=true,
  pdflang={zh-CN},
  pdftitle={无穷范畴导论（中文译本）},
  pdfauthor={Markus Land（原著）},
}

\usepackage{imakeidx}
\makeindex[title={译本术语索引},intoc,noautomatic]
\pdfstringdefDisableCommands{\def\index#1{}}
\makeatletter
\let\ic@tableofcontents\tableofcontents
\def\tableofcontents{\begingroup\def\index##1{}\ic@tableofcontents\endgroup}
\makeatother

\usepackage[capitalise,noabbrev]{cleveref}
\crefname{theorem}     {\trhTheorem}     {\trhTheorem}
\crefname{proposition} {\trhProposition} {\trhProposition}
\crefname{lemma}       {\trhLemma}       {\trhLemma}
\crefname{corollary}   {\trhCorollary}   {\trhCorollary}
\crefname{fact}        {\trhFact}        {\trhFact}
\crefname{definition}  {\trhDefinition}  {\trhDefinition}
\crefname{example}     {\trhExample}     {\trhExample}
\crefname{notation}    {\trhNotation}    {\trhNotation}
\crefname{convention}  {\trhConvention}  {\trhConvention}
\crefname{construction}{\trhConstruction}{\trhConstruction}
\crefname{remark}      {\trhRemark}      {\trhRemark}
\crefname{observation} {\trhObservation} {\trhObservation}
\crefname{warning}     {\trhWarning}     {\trhWarning}
\crefname{exercise}    {\trhExercise}    {\trhExercise}
\crefname{figure}      {图}{图}
\crefname{table}       {表}{表}
\crefname{equation}    {}{}  %
\crefformat{equation}{#2(#1)#3}

\usepackage{booktabs,longtable,multicol}
\newcommand{\BookGlossEntry}[2]{%
  \par\noindent\hangindent=1em\hangafter=1 #1\quad #2\par}
\newcommand{\BookSourceIndexEntry}[2]{%
  \par\begingroup\leftskip=#1em\relax\hangindent=1em\relax
  \noindent #2\par\endgroup}
\makeatletter
\newif\ifBookIndexing
\BookIndexingfalse
\newcommand{\BookBoldPage}[1]{\textbf{\hyperpage{#1}}}
\DeclareRobustCommand{\BookRecordIndex}[1]{%
  \ifBookIndexing
    \expandafter\index\expandafter{\csname BookIndex@#1\endcsname|hyperpage}%
  \fi}
\DeclareRobustCommand{\BookIndexDefinition}[1]{%
  \ifBookIndexing
    \expandafter\index\expandafter{\csname BookIndex@#1\endcsname|BookBoldPage}%
  \fi}
\DeclareRobustCommand{\BookTerm}[2]{%
  \BookRecordIndex{#1}%
  \begingroup\let\BookRecordIndex\@gobble #2\endgroup}
\AddToHook{cmd/@outputpage/before}{\let\BookRecordIndex\@gobble}
\pdfstringdefDisableCommands{%
  \def\BookTerm#1#2{#2}%
  \let\BookRecordIndex\@gobble
  \let\BookIndexDefinition\@gobble}
\expandafter\def\csname BookIndex@T0001\endcsname{fanchou@∞-范畴（\ensuremath{\infty}-category）}
\expandafter\let\csname BookOriginal@trInfinityCategory\endcsname\trInfinityCategory
\renewcommand{\trInfinityCategory}{\BookTerm{T0001}{\csname BookOriginal@trInfinityCategory\endcsname}}
\expandafter\def\csname BookIndex@T0002\endcsname{nifanchou@拟范畴（quasi-category）}
\expandafter\let\csname BookOriginal@trQuasiCategory\endcsname\trQuasiCategory
\renewcommand{\trQuasiCategory}{\BookTerm{T0002}{\csname BookOriginal@trQuasiCategory\endcsname}}
\expandafter\def\csname BookIndex@T0003\endcsname{fanchou@范畴（category）}
\expandafter\let\csname BookOriginal@trCategory\endcsname\trCategory
\renewcommand{\trCategory}{\BookTerm{T0003}{\csname BookOriginal@trCategory\endcsname}}
\expandafter\def\csname BookIndex@T0004\endcsname{danchunji@单纯集（simplicial set）}
\expandafter\let\csname BookOriginal@trSimplicialSet\endcsname\trSimplicialSet
\renewcommand{\trSimplicialSet}{\BookTerm{T0004}{\csname BookOriginal@trSimplicialSet\endcsname}}
\expandafter\def\csname BookIndex@T0005\endcsname{danchunfanchou@单纯范畴（simplicial category）}
\expandafter\let\csname BookOriginal@trSimplicialCategory\endcsname\trSimplicialCategory
\renewcommand{\trSimplicialCategory}{\BookTerm{T0005}{\csname BookOriginal@trSimplicialCategory\endcsname}}
\expandafter\def\csname BookIndex@T0006\endcsname{shenjing@神经（nerve）}
\expandafter\let\csname BookOriginal@trNerve\endcsname\trNerve
\renewcommand{\trNerve}{\BookTerm{T0006}{\csname BookOriginal@trNerve\endcsname}}
\expandafter\def\csname BookIndex@T0007\endcsname{jiao@角（horn）}
\expandafter\let\csname BookOriginal@trHorn\endcsname\trHorn
\renewcommand{\trHorn}{\BookTerm{T0007}{\csname BookOriginal@trHorn\endcsname}}
\expandafter\def\csname BookIndex@T0008\endcsname{danxing@单形（simplex）}
\expandafter\let\csname BookOriginal@trSimplex\endcsname\trSimplex
\renewcommand{\trSimplex}{\BookTerm{T0008}{\csname BookOriginal@trSimplex\endcsname}}
\expandafter\def\csname BookIndex@T0009\endcsname{kanfuxing@Kan 复形（Kan complex）}
\expandafter\let\csname BookOriginal@trKanComplex\endcsname\trKanComplex
\renewcommand{\trKanComplex}{\BookTerm{T0009}{\csname BookOriginal@trKanComplex\endcsname}}
\expandafter\def\csname BookIndex@T0010\endcsname{kanxianweihua@Kan 纤维化（Kan fibration）}
\expandafter\let\csname BookOriginal@trKanFibration\endcsname\trKanFibration
\renewcommand{\trKanFibration}{\BookTerm{T0010}{\csname BookOriginal@trKanFibration\endcsname}}
\expandafter\def\csname BookIndex@T0011\endcsname{ruodengjia@弱等价（weak equivalence）}
\expandafter\let\csname BookOriginal@trWeakEquivalence\endcsname\trWeakEquivalence
\renewcommand{\trWeakEquivalence}{\BookTerm{T0011}{\csname BookOriginal@trWeakEquivalence\endcsname}}
\expandafter\def\csname BookIndex@T0012\endcsname{ruoyaerdengjia@若亚尔等价（Joyal equivalence）}
\expandafter\let\csname BookOriginal@trJoyalEquivalence\endcsname\trJoyalEquivalence
\renewcommand{\trJoyalEquivalence}{\BookTerm{T0012}{\csname BookOriginal@trJoyalEquivalence\endcsname}}
\expandafter\def\csname BookIndex@T0013\endcsname{yingshekongjian@映射空间（mapping space）}
\expandafter\let\csname BookOriginal@trMappingSpace\endcsname\trMappingSpace
\renewcommand{\trMappingSpace}{\BookTerm{T0013}{\csname BookOriginal@trMappingSpace\endcsname}}
\expandafter\def\csname BookIndex@T0014\endcsname{hanzi@函子（functor）}
\expandafter\let\csname BookOriginal@trFunctor\endcsname\trFunctor
\renewcommand{\trFunctor}{\BookTerm{T0014}{\csname BookOriginal@trFunctor\endcsname}}
\expandafter\def\csname BookIndex@T0015\endcsname{ziranbianhuan@自然变换（natural transformation）}
\expandafter\let\csname BookOriginal@trNaturalTransformation\endcsname\trNaturalTransformation
\renewcommand{\trNaturalTransformation}{\BookTerm{T0015}{\csname BookOriginal@trNaturalTransformation\endcsname}}
\expandafter\def\csname BookIndex@T0016\endcsname{dengjia@等价（equivalence）}
\expandafter\let\csname BookOriginal@trEquivalence\endcsname\trEquivalence
\renewcommand{\trEquivalence}{\BookTerm{T0016}{\csname BookOriginal@trEquivalence\endcsname}}
\expandafter\def\csname BookIndex@T0017\endcsname{jixian@极限（limit）}
\expandafter\let\csname BookOriginal@trLimit\endcsname\trLimit
\renewcommand{\trLimit}{\BookTerm{T0017}{\csname BookOriginal@trLimit\endcsname}}
\expandafter\def\csname BookIndex@T0018\endcsname{yujixian@余极限（colimit）}
\expandafter\let\csname BookOriginal@trColimit\endcsname\trColimit
\renewcommand{\trColimit}{\BookTerm{T0018}{\csname BookOriginal@trColimit\endcsname}}
\expandafter\def\csname BookIndex@T0019\endcsname{zhongduixiang@终对象（terminal object）}
\expandafter\let\csname BookOriginal@trTerminalObject\endcsname\trTerminalObject
\renewcommand{\trTerminalObject}{\BookTerm{T0019}{\csname BookOriginal@trTerminalObject\endcsname}}
\expandafter\def\csname BookIndex@T0020\endcsname{shiduixiang@始对象（initial object）}
\expandafter\let\csname BookOriginal@trInitialObject\endcsname\trInitialObject
\renewcommand{\trInitialObject}{\BookTerm{T0020}{\csname BookOriginal@trInitialObject\endcsname}}
\expandafter\def\csname BookIndex@T0021\endcsname{bansui@伴随（adjunction）}
\expandafter\let\csname BookOriginal@trAdjunction\endcsname\trAdjunction
\renewcommand{\trAdjunction}{\BookTerm{T0021}{\csname BookOriginal@trAdjunction\endcsname}}
\expandafter\def\csname BookIndex@T0022\endcsname{zuobansui@左伴随（left adjoint）}
\expandafter\let\csname BookOriginal@trLeftAdjoint\endcsname\trLeftAdjoint
\renewcommand{\trLeftAdjoint}{\BookTerm{T0022}{\csname BookOriginal@trLeftAdjoint\endcsname}}
\expandafter\def\csname BookIndex@T0023\endcsname{youbansui@右伴随（right adjoint）}
\expandafter\let\csname BookOriginal@trRightAdjoint\endcsname\trRightAdjoint
\renewcommand{\trRightAdjoint}{\BookTerm{T0023}{\csname BookOriginal@trRightAdjoint\endcsname}}
\expandafter\def\csname BookIndex@T0024\endcsname{bansuihanzidingli@伴随函子定理（adjoint functor theorem）}
\expandafter\let\csname BookOriginal@trAdjointFunctorTheorem\endcsname\trAdjointFunctorTheorem
\renewcommand{\trAdjointFunctorTheorem}{\BookTerm{T0024}{\csname BookOriginal@trAdjointFunctorTheorem\endcsname}}
\expandafter\def\csname BookIndex@T0025\endcsname{jubuhua@局部化（localization）}
\expandafter\let\csname BookOriginal@trLocalization\endcsname\trLocalization
\renewcommand{\trLocalization}{\BookTerm{T0025}{\csname BookOriginal@trLocalization\endcsname}}
\expandafter\def\csname BookIndex@T0026\endcsname{xianweihua@纤维化（fibration）}
\expandafter\let\csname BookOriginal@trFibration\endcsname\trFibration
\renewcommand{\trFibration}{\BookTerm{T0026}{\csname BookOriginal@trFibration\endcsname}}
\expandafter\def\csname BookIndex@T0027\endcsname{neixianweihua@内纤维化（inner fibration）}
\expandafter\let\csname BookOriginal@trInnerFibration\endcsname\trInnerFibration
\renewcommand{\trInnerFibration}{\BookTerm{T0027}{\csname BookOriginal@trInnerFibration\endcsname}}
\expandafter\def\csname BookIndex@T0028\endcsname{neilingdiao@内零调（inner anodyne）}
\expandafter\let\csname BookOriginal@trInnerAnodyne\endcsname\trInnerAnodyne
\renewcommand{\trInnerAnodyne}{\BookTerm{T0028}{\csname BookOriginal@trInnerAnodyne\endcsname}}
\expandafter\def\csname BookIndex@T0029\endcsname{lingdiaoyingshe@零调映射（anodyne map）}
\expandafter\let\csname BookOriginal@trAnodyneMap\endcsname\trAnodyneMap
\renewcommand{\trAnodyneMap}{\BookTerm{T0029}{\csname BookOriginal@trAnodyneMap\endcsname}}
\expandafter\def\csname BookIndex@T0030\endcsname{dikaerxianweihua@笛卡尔纤维化（cartesian fibration）}
\expandafter\let\csname BookOriginal@trCartesianFibration\endcsname\trCartesianFibration
\renewcommand{\trCartesianFibration}{\BookTerm{T0030}{\csname BookOriginal@trCartesianFibration\endcsname}}
\expandafter\def\csname BookIndex@T0031\endcsname{yudikaerxianweihua@余笛卡尔纤维化（cocartesian fibration）}
\expandafter\let\csname BookOriginal@trCocartesianFibration\endcsname\trCocartesianFibration
\renewcommand{\trCocartesianFibration}{\BookTerm{T0031}{\csname BookOriginal@trCocartesianFibration\endcsname}}
\expandafter\def\csname BookIndex@T0032\endcsname{daibiaojidanchunji@带标记单纯集（marked simplicial set）}
\expandafter\let\csname BookOriginal@trMarkedSimplicialSet\endcsname\trMarkedSimplicialSet
\renewcommand{\trMarkedSimplicialSet}{\BookTerm{T0032}{\csname BookOriginal@trMarkedSimplicialSet\endcsname}}
\expandafter\def\csname BookIndex@T0033\endcsname{daibiaojilingdiao@带标记零调（marked anodyne）}
\expandafter\let\csname BookOriginal@trMarkedAnodyne\endcsname\trMarkedAnodyne
\renewcommand{\trMarkedAnodyne}{\BookTerm{T0033}{\csname BookOriginal@trMarkedAnodyne\endcsname}}
\expandafter\def\csname BookIndex@T0034\endcsname{lazhi@拉直（straightening）}
\expandafter\let\csname BookOriginal@trStraightening\endcsname\trStraightening
\renewcommand{\trStraightening}{\BookTerm{T0034}{\csname BookOriginal@trStraightening\endcsname}}
\expandafter\def\csname BookIndex@T0035\endcsname{fanlazhi@反拉直（unstraightening）}
\expandafter\let\csname BookOriginal@trUnstraightening\endcsname\trUnstraightening
\renewcommand{\trUnstraightening}{\BookTerm{T0035}{\csname BookOriginal@trUnstraightening\endcsname}}
\expandafter\def\csname BookIndex@T0036\endcsname{niujianfanchou@扭箭范畴（twisted arrow category）}
\expandafter\let\csname BookOriginal@trTwistedArrowCategory\endcsname\trTwistedArrowCategory
\renewcommand{\trTwistedArrowCategory}{\BookTerm{T0036}{\csname BookOriginal@trTwistedArrowCategory\endcsname}}
\expandafter\def\csname BookIndex@T0037\endcsname{lian@联（join）}
\expandafter\let\csname BookOriginal@trJoin\endcsname\trJoin
\renewcommand{\trJoin}{\BookTerm{T0037}{\csname BookOriginal@trJoin\endcsname}}
\expandafter\def\csname BookIndex@T0038\endcsname{panglian@胖联（fat join）}
\expandafter\let\csname BookOriginal@trFatJoin\endcsname\trFatJoin
\renewcommand{\trFatJoin}{\BookTerm{T0038}{\csname BookOriginal@trFatJoin\endcsname}}
\expandafter\def\csname BookIndex@T0039\endcsname{pangqiepian@胖切片（fat slice）}
\expandafter\let\csname BookOriginal@trFatSlice\endcsname\trFatSlice
\renewcommand{\trFatSlice}{\BookTerm{T0039}{\csname BookOriginal@trFatSlice\endcsname}}
\expandafter\def\csname BookIndex@T0040\endcsname{tonglundengjia@同伦等价（homotopy equivalence）}
\expandafter\let\csname BookOriginal@trHomotopyEquivalence\endcsname\trHomotopyEquivalence
\renewcommand{\trHomotopyEquivalence}{\BookTerm{T0040}{\csname BookOriginal@trHomotopyEquivalence\endcsname}}
\expandafter\def\csname BookIndex@T0041\endcsname{qiepianfanchou@切片范畴（slice category）}
\expandafter\let\csname BookOriginal@trSliceCategory\endcsname\trSliceCategory
\renewcommand{\trSliceCategory}{\BookTerm{T0041}{\csname BookOriginal@trSliceCategory\endcsname}}
\expandafter\def\csname BookIndex@T0042\endcsname{gongweide@共尾的（cofinal）}
\expandafter\let\csname BookOriginal@trCofinal\endcsname\trCofinal
\renewcommand{\trCofinal}{\BookTerm{T0042}{\csname BookOriginal@trCofinal\endcsname}}
\expandafter\def\csname BookIndex@T0043\endcsname{gongshide@共始的（coinitial）}
\expandafter\let\csname BookOriginal@trCoinitial\endcsname\trCoinitial
\renewcommand{\trCoinitial}{\BookTerm{T0043}{\csname BookOriginal@trCoinitial\endcsname}}
\expandafter\def\csname BookIndex@T0044\endcsname{quanzhongshide@全忠实的（fully faithful）}
\expandafter\let\csname BookOriginal@trFullyFaithful\endcsname\trFullyFaithful
\renewcommand{\trFullyFaithful}{\BookTerm{T0044}{\csname BookOriginal@trFullyFaithful\endcsname}}
\expandafter\def\csname BookIndex@T0045\endcsname{benzhimande@本质满的（essentially surjective）}
\expandafter\let\csname BookOriginal@trEssentiallySurjective\endcsname\trEssentiallySurjective
\renewcommand{\trEssentiallySurjective}{\BookTerm{T0045}{\csname BookOriginal@trEssentiallySurjective\endcsname}}
\expandafter\def\csname BookIndex@T0046\endcsname{tonglun@同伦（homotopy）}
\expandafter\let\csname BookOriginal@trHomotopy\endcsname\trHomotopy
\renewcommand{\trHomotopy}{\BookTerm{T0046}{\csname BookOriginal@trHomotopy\endcsname}}
\expandafter\def\csname BookIndex@T0047\endcsname{luguode@滤过的（filtered）}
\expandafter\let\csname BookOriginal@trFiltered\endcsname\trFiltered
\renewcommand{\trFiltered}{\BookTerm{T0047}{\csname BookOriginal@trFiltered\endcsname}}
\expandafter\def\csname BookIndex@T0048\endcsname{fanchoulun@范畴论（category theory）}
\expandafter\let\csname BookOriginal@trCategoryTheory\endcsname\trCategoryTheory
\renewcommand{\trCategoryTheory}{\BookTerm{T0048}{\csname BookOriginal@trCategoryTheory\endcsname}}
\expandafter\def\csname BookIndex@T0049\endcsname{gaojiefanchoulun@高阶范畴论（higher category theory）}
\expandafter\let\csname BookOriginal@trHigherCategoryTheory\endcsname\trHigherCategoryTheory
\renewcommand{\trHigherCategoryTheory}{\BookTerm{T0049}{\csname BookOriginal@trHigherCategoryTheory\endcsname}}
\expandafter\def\csname BookIndex@T0050\endcsname{tonglunlun@同伦论（homotopy theory）}
\expandafter\let\csname BookOriginal@trHomotopyTheory\endcsname\trHomotopyTheory
\renewcommand{\trHomotopyTheory}{\BookTerm{T0050}{\csname BookOriginal@trHomotopyTheory\endcsname}}
\expandafter\def\csname BookIndex@T0051\endcsname{moxingfanchou@模型范畴（model category）}
\expandafter\let\csname BookOriginal@trModelCategory\endcsname\trModelCategory
\renewcommand{\trModelCategory}{\BookTerm{T0051}{\csname BookOriginal@trModelCategory\endcsname}}
\expandafter\def\csname BookIndex@T0052\endcsname{daishutapu@代数拓扑（algebraic topology）}
\expandafter\let\csname BookOriginal@trAlgebraicTopology\endcsname\trAlgebraicTopology
\renewcommand{\trAlgebraicTopology}{\BookTerm{T0052}{\csname BookOriginal@trAlgebraicTopology\endcsname}}
\expandafter\def\csname BookIndex@T0053\endcsname{tapukongjian@拓扑空间（topological space）}
\expandafter\let\csname BookOriginal@trTopologicalSpace\endcsname\trTopologicalSpace
\renewcommand{\trTopologicalSpace}{\BookTerm{T0053}{\csname BookOriginal@trTopologicalSpace\endcsname}}
\expandafter\def\csname BookIndex@T0054\endcsname{chongshi@充实（enrichment）}
\expandafter\let\csname BookOriginal@trEnrichment\endcsname\trEnrichment
\renewcommand{\trEnrichment}{\BookTerm{T0054}{\csname BookOriginal@trEnrichment\endcsname}}
\expandafter\def\csname BookIndex@T0055\endcsname{ningjushenjing@凝聚神经（coherent nerve）}
\expandafter\let\csname BookOriginal@trCoherentNerve\endcsname\trCoherentNerve
\renewcommand{\trCoherentNerve}{\BookTerm{T0055}{\csname BookOriginal@trCoherentNerve\endcsname}}
\expandafter\def\csname BookIndex@T0056\endcsname{duixiang@对象（object）}
\expandafter\let\csname BookOriginal@trObject\endcsname\trObject
\renewcommand{\trObject}{\BookTerm{T0056}{\csname BookOriginal@trObject\endcsname}}
\expandafter\def\csname BookIndex@T0057\endcsname{taishe@态射（morphism）}
\expandafter\let\csname BookOriginal@trMorphism\endcsname\trMorphism
\renewcommand{\trMorphism}{\BookTerm{T0057}{\csname BookOriginal@trMorphism\endcsname}}
\expandafter\def\csname BookIndex@T0058\endcsname{dantaishe@单态射（monomorphism）}
\expandafter\let\csname BookOriginal@trMonomorphism\endcsname\trMonomorphism
\renewcommand{\trMonomorphism}{\BookTerm{T0058}{\csname BookOriginal@trMonomorphism\endcsname}}
\expandafter\def\csname BookIndex@T0059\endcsname{yuceng@预层（presheaf）}
\expandafter\let\csname BookOriginal@trPresheaf\endcsname\trPresheaf
\renewcommand{\trPresheaf}{\BookTerm{T0059}{\csname BookOriginal@trPresheaf\endcsname}}
\expandafter\def\csname BookIndex@T0060\endcsname{qunpei@群胚（groupoid）}
\expandafter\let\csname BookOriginal@trGroupoid\endcsname\trGroupoid
\renewcommand{\trGroupoid}{\BookTerm{T0060}{\csname BookOriginal@trGroupoid\endcsname}}
\expandafter\def\csname BookIndex@T0061\endcsname{tuichu@推出（pushout）}
\expandafter\let\csname BookOriginal@trPushout\endcsname\trPushout
\renewcommand{\trPushout}{\BookTerm{T0061}{\csname BookOriginal@trPushout\endcsname}}
\expandafter\def\csname BookIndex@T0062\endcsname{lahui@拉回（pullback）}
\expandafter\let\csname BookOriginal@trPullback\endcsname\trPullback
\renewcommand{\trPullback}{\BookTerm{T0062}{\csname BookOriginal@trPullback\endcsname}}
\expandafter\def\csname BookIndex@T0063\endcsname{kankuozhang@Kan 扩张（Kan extension）}
\expandafter\let\csname BookOriginal@trKanExtension\endcsname\trKanExtension
\renewcommand{\trKanExtension}{\BookTerm{T0063}{\csname BookOriginal@trKanExtension\endcsname}}
\expandafter\def\csname BookIndex@T0064\endcsname{ji@脊（spine）}
\expandafter\let\csname BookOriginal@trSpine\endcsname\trSpine
\renewcommand{\trSpine}{\BookTerm{T0064}{\csname BookOriginal@trSpine\endcsname}}
\expandafter\def\csname BookIndex@T0065\endcsname{gujia@骨架（skeleton）}
\expandafter\let\csname BookOriginal@trSkeleton\endcsname\trSkeleton
\renewcommand{\trSkeleton}{\BookTerm{T0065}{\csname BookOriginal@trSkeleton\endcsname}}
\expandafter\def\csname BookIndex@T0066\endcsname{yugujia@余骨架（coskeleton）}
\expandafter\let\csname BookOriginal@trCoskeleton\endcsname\trCoskeleton
\renewcommand{\trCoskeleton}{\BookTerm{T0066}{\csname BookOriginal@trCoskeleton\endcsname}}
\expandafter\def\csname BookIndex@T0067\endcsname{jiheshixian@几何实现（geometric realization）}
\expandafter\let\csname BookOriginal@trGeometricRealization\endcsname\trGeometricRealization
\renewcommand{\trGeometricRealization}{\BookTerm{T0067}{\csname BookOriginal@trGeometricRealization\endcsname}}
\expandafter\def\csname BookIndex@T0068\endcsname{qiyidanchunji@奇异单纯集（singular simplicial set）}
\expandafter\let\csname BookOriginal@trSingularSimplicialSet\endcsname\trSingularSimplicialSet
\renewcommand{\trSingularSimplicialSet}{\BookTerm{T0068}{\csname BookOriginal@trSingularSimplicialSet\endcsname}}
\expandafter\def\csname BookIndex@T0069\endcsname{yudanchun@余单纯（cosimplicial）}
\expandafter\let\csname BookOriginal@trCosimplicial\endcsname\trCosimplicial
\renewcommand{\trCosimplicial}{\BookTerm{T0069}{\csname BookOriginal@trCosimplicial\endcsname}}
\expandafter\def\csname BookIndex@T0070\endcsname{mianyingshe@面映射（face map）}
\expandafter\let\csname BookOriginal@trFaceMap\endcsname\trFaceMap
\renewcommand{\trFaceMap}{\BookTerm{T0070}{\csname BookOriginal@trFaceMap\endcsname}}
\expandafter\def\csname BookIndex@T0071\endcsname{tuihuayingshe@退化映射（degeneracy map）}
\expandafter\let\csname BookOriginal@trDegeneracyMap\endcsname\trDegeneracyMap
\renewcommand{\trDegeneracyMap}{\BookTerm{T0071}{\csname BookOriginal@trDegeneracyMap\endcsname}}
\expandafter\def\csname BookIndex@T0072\endcsname{shuangwanbei@双完备（bicomplete）}
\expandafter\let\csname BookOriginal@trBicomplete\endcsname\trBicomplete
\renewcommand{\trBicomplete}{\BookTerm{T0072}{\csname BookOriginal@trBicomplete\endcsname}}
\expandafter\def\csname BookIndex@T0073\endcsname{yuzhui@余锥（cocone）}
\expandafter\let\csname BookOriginal@trCocone\endcsname\trCocone
\renewcommand{\trCocone}{\BookTerm{T0073}{\csname BookOriginal@trCocone\endcsname}}
\expandafter\def\csname BookIndex@T0074\endcsname{tishengwenti@提升问题（lifting problem）}
\expandafter\let\csname BookOriginal@trLiftingProblem\endcsname\trLiftingProblem
\renewcommand{\trLiftingProblem}{\BookTerm{T0074}{\csname BookOriginal@trLiftingProblem\endcsname}}
\expandafter\def\csname BookIndex@T0075\endcsname{hechengqi@合成器（composer）}
\expandafter\let\csname BookOriginal@trComposer\endcsname\trComposer
\renewcommand{\trComposer}{\BookTerm{T0075}{\csname BookOriginal@trComposer\endcsname}}
\expandafter\def\csname BookIndex@T0076\endcsname{yaobanjiegou@幺半结构（monoidal structure）}
\expandafter\let\csname BookOriginal@trMonoidalStructure\endcsname\trMonoidalStructure
\renewcommand{\trMonoidalStructure}{\BookTerm{T0076}{\csname BookOriginal@trMonoidalStructure\endcsname}}
\expandafter\def\csname BookIndex@T0077\endcsname{yaobanji@幺半积（monoidal product）}
\expandafter\let\csname BookOriginal@trMonoidalProduct\endcsname\trMonoidalProduct
\renewcommand{\trMonoidalProduct}{\BookTerm{T0077}{\csname BookOriginal@trMonoidalProduct\endcsname}}
\expandafter\def\csname BookIndex@T0078\endcsname{yaobanbansui@幺半伴随（monoidal adjunction）}
\expandafter\let\csname BookOriginal@trMonoidalAdjunction\endcsname\trMonoidalAdjunction
\renewcommand{\trMonoidalAdjunction}{\BookTerm{T0078}{\csname BookOriginal@trMonoidalAdjunction\endcsname}}
\expandafter\def\csname BookIndex@T0079\endcsname{songyaobanhanzi@松幺半函子（lax monoidal functor）}
\expandafter\let\csname BookOriginal@trLaxMonoidalFunctor\endcsname\trLaxMonoidalFunctor
\renewcommand{\trLaxMonoidalFunctor}{\BookTerm{T0079}{\csname BookOriginal@trLaxMonoidalFunctor\endcsname}}
\expandafter\def\csname BookIndex@T0080\endcsname{fansongyaobanhanzi@反松幺半函子（oplax monoidal functor）}
\expandafter\let\csname BookOriginal@trOplaxMonoidalFunctor\endcsname\trOplaxMonoidalFunctor
\renewcommand{\trOplaxMonoidalFunctor}{\BookTerm{T0080}{\csname BookOriginal@trOplaxMonoidalFunctor\endcsname}}
\expandafter\def\csname BookIndex@T0081\endcsname{difanchou@底范畴（underlying category）}
\expandafter\let\csname BookOriginal@trUnderlyingCategory\endcsname\trUnderlyingCategory
\renewcommand{\trUnderlyingCategory}{\BookTerm{T0081}{\csname BookOriginal@trUnderlyingCategory\endcsname}}
\expandafter\def\csname BookIndex@T0082\endcsname{neimian@内面（inner face）}
\expandafter\let\csname BookOriginal@trInnerFace\endcsname\trInnerFace
\renewcommand{\trInnerFace}{\BookTerm{T0082}{\csname BookOriginal@trInnerFace\endcsname}}
\expandafter\def\csname BookIndex@T0083\endcsname{chuankonglifangti@穿孔立方体（punctured cube）}
\expandafter\let\csname BookOriginal@trPuncturedCube\endcsname\trPuncturedCube
\renewcommand{\trPuncturedCube}{\BookTerm{T0083}{\csname BookOriginal@trPuncturedCube\endcsname}}
\expandafter\def\csname BookIndex@T0084\endcsname{2hanzi@2-函子（2-functor）}
\expandafter\let\csname BookOriginal@trTwoFunctor\endcsname\trTwoFunctor
\renewcommand{\trTwoFunctor}{\BookTerm{T0084}{\csname BookOriginal@trTwoFunctor\endcsname}}
\expandafter\def\csname BookIndex@T0085\endcsname{2taishe@2-态射（2-morphism）}
\expandafter\let\csname BookOriginal@trTwoMorphism\endcsname\trTwoMorphism
\renewcommand{\trTwoMorphism}{\BookTerm{T0085}{\csname BookOriginal@trTwoMorphism\endcsname}}
\expandafter\def\csname BookIndex@T0086\endcsname{cwfuxing@CW 复形（CW complex）}
\expandafter\let\csname BookOriginal@trCWComplex\endcsname\trCWComplex
\renewcommand{\trCWComplex}{\BookTerm{T0086}{\csname BookOriginal@trCWComplex\endcsname}}
\expandafter\def\csname BookIndex@T0087\endcsname{ruoquanzhongshide@弱全忠实的（weakly fully faithful）}
\expandafter\let\csname BookOriginal@trWeaklyFullyFaithful\endcsname\trWeaklyFullyFaithful
\renewcommand{\trWeaklyFullyFaithful}{\BookTerm{T0087}{\csname BookOriginal@trWeaklyFullyFaithful\endcsname}}
\expandafter\def\csname BookIndex@T0088\endcsname{ruobenzhimande@弱本质满的（weakly essentially surjective）}
\expandafter\let\csname BookOriginal@trWeaklyEssentiallySurjective\endcsname\trWeaklyEssentiallySurjective
\renewcommand{\trWeaklyEssentiallySurjective}{\BookTerm{T0088}{\csname BookOriginal@trWeaklyEssentiallySurjective\endcsname}}
\expandafter\def\csname BookIndex@T0089\endcsname{zuoxianweihua@左纤维化（left fibration）}
\expandafter\let\csname BookOriginal@trLeftFibration\endcsname\trLeftFibration
\renewcommand{\trLeftFibration}{\BookTerm{T0089}{\csname BookOriginal@trLeftFibration\endcsname}}
\expandafter\def\csname BookIndex@T0090\endcsname{youxianweihua@右纤维化（right fibration）}
\expandafter\let\csname BookOriginal@trRightFibration\endcsname\trRightFibration
\renewcommand{\trRightFibration}{\BookTerm{T0090}{\csname BookOriginal@trRightFibration\endcsname}}
\expandafter\def\csname BookIndex@T0091\endcsname{pingfanxianweihua@平凡纤维化（trivial fibration）}
\expandafter\let\csname BookOriginal@trTrivialFibration\endcsname\trTrivialFibration
\renewcommand{\trTrivialFibration}{\BookTerm{T0091}{\csname BookOriginal@trTrivialFibration\endcsname}}
\expandafter\def\csname BookIndex@T0092\endcsname{lingdiao@零调（anodyne）}
\expandafter\let\csname BookOriginal@trAnodyne\endcsname\trAnodyne
\renewcommand{\trAnodyne}{\BookTerm{T0092}{\csname BookOriginal@trAnodyne\endcsname}}
\expandafter\def\csname BookIndex@T0093\endcsname{zuolingdiao@左零调（left anodyne）}
\expandafter\let\csname BookOriginal@trLeftAnodyne\endcsname\trLeftAnodyne
\renewcommand{\trLeftAnodyne}{\BookTerm{T0093}{\csname BookOriginal@trLeftAnodyne\endcsname}}
\expandafter\def\csname BookIndex@T0094\endcsname{youlingdiao@右零调（right anodyne）}
\expandafter\let\csname BookOriginal@trRightAnodyne\endcsname\trRightAnodyne
\renewcommand{\trRightAnodyne}{\BookTerm{T0094}{\csname BookOriginal@trRightAnodyne\endcsname}}
\expandafter\def\csname BookIndex@T0095\endcsname{baohede@饱和的（saturated）}
\expandafter\let\csname BookOriginal@trSaturated\endcsname\trSaturated
\renewcommand{\trSaturated}{\BookTerm{T0095}{\csname BookOriginal@trSaturated\endcsname}}
\expandafter\def\csname BookIndex@T0096\endcsname{baoheji@饱和集（saturated set）}
\expandafter\let\csname BookOriginal@trSaturatedSet\endcsname\trSaturatedSet
\renewcommand{\trSaturatedSet}{\BookTerm{T0096}{\csname BookOriginal@trSaturatedSet\endcsname}}
\expandafter\def\csname BookIndex@T0097\endcsname{baohebibao@饱和闭包（saturation）}
\expandafter\let\csname BookOriginal@trSaturation\endcsname\trSaturation
\renewcommand{\trSaturation}{\BookTerm{T0097}{\csname BookOriginal@trSaturation\endcsname}}
\expandafter\def\csname BookIndex@T0098\endcsname{shousuo@收缩（retract）}
\expandafter\let\csname BookOriginal@trRetract\endcsname\trRetract
\renewcommand{\trRetract}{\BookTerm{T0098}{\csname BookOriginal@trRetract\endcsname}}
\expandafter\def\csname BookIndex@T0099\endcsname{youtishengxingzhi@右提升性质（right lifting property）}
\expandafter\let\csname BookOriginal@trRightLiftingProperty\endcsname\trRightLiftingProperty
\renewcommand{\trRightLiftingProperty}{\BookTerm{T0099}{\csname BookOriginal@trRightLiftingProperty\endcsname}}
\expandafter\def\csname BookIndex@T0100\endcsname{zuotishengxingzhi@左提升性质（left lifting property）}
\expandafter\let\csname BookOriginal@trLeftLiftingProperty\endcsname\trLeftLiftingProperty
\renewcommand{\trLeftLiftingProperty}{\BookTerm{T0100}{\csname BookOriginal@trLeftLiftingProperty\endcsname}}
\expandafter\def\csname BookIndex@T0101\endcsname{xiaoduixianglunzheng@小对象论证（small object argument）}
\expandafter\let\csname BookOriginal@trSmallObjectArgument\endcsname\trSmallObjectArgument
\renewcommand{\trSmallObjectArgument}{\BookTerm{T0101}{\csname BookOriginal@trSmallObjectArgument\endcsname}}
\expandafter\def\csname BookIndex@T0102\endcsname{jiaobaohan@角包含（horn inclusion）}
\expandafter\let\csname BookOriginal@trHornInclusion\endcsname\trHornInclusion
\renewcommand{\trHornInclusion}{\BookTerm{T0102}{\csname BookOriginal@trHornInclusion\endcsname}}
\expandafter\def\csname BookIndex@T0103\endcsname{bianjiebaohan@边界包含（boundary inclusion）}
\expandafter\let\csname BookOriginal@trBoundaryInclusion\endcsname\trBoundaryInclusion
\renewcommand{\trBoundaryInclusion}{\BookTerm{T0103}{\csname BookOriginal@trBoundaryInclusion\endcsname}}
\expandafter\def\csname BookIndex@T0104\endcsname{jibaohan@脊包含（spine inclusion）}
\expandafter\let\csname BookOriginal@trSpineInclusion\endcsname\trSpineInclusion
\renewcommand{\trSpineInclusion}{\BookTerm{T0104}{\csname BookOriginal@trSpineInclusion\endcsname}}
\expandafter\def\csname BookIndex@T0105\endcsname{jinde@紧的（compact）}
\expandafter\let\csname BookOriginal@trCompact\endcsname\trCompact
\renewcommand{\trCompact}{\BookTerm{T0105}{\csname BookOriginal@trCompact\endcsname}}
\expandafter\def\csname BookIndex@T0106\endcsname{jiemian@截面（section）}
\expandafter\let\csname BookOriginal@trSectionMap\endcsname\trSectionMap
\renewcommand{\trSectionMap}{\BookTerm{T0106}{\csname BookOriginal@trSectionMap\endcsname}}
\expandafter\def\csname BookIndex@T0107\endcsname{tuichuji@推出积（pushout product）}
\expandafter\let\csname BookOriginal@trPushoutProduct\endcsname\trPushoutProduct
\renewcommand{\trPushoutProduct}{\BookTerm{T0107}{\csname BookOriginal@trPushoutProduct\endcsname}}
\expandafter\def\csname BookIndex@T0108\endcsname{ruoyaerxianweihua@若亚尔纤维化（Joyal fibration）}
\expandafter\let\csname BookOriginal@trJoyalFibration\endcsname\trJoyalFibration
\renewcommand{\trJoyalFibration}{\BookTerm{T0108}{\csname BookOriginal@trJoyalFibration\endcsname}}
\expandafter\def\csname BookIndex@T0109\endcsname{moxingjiegou@模型结构（model structure）}
\expandafter\let\csname BookOriginal@trModelStructure\endcsname\trModelStructure
\renewcommand{\trModelStructure}{\BookTerm{T0109}{\csname BookOriginal@trModelStructure\endcsname}}
\expandafter\def\csname BookIndex@T0110\endcsname{fenleiqi@分类器（classifier）}
\expandafter\let\csname BookOriginal@trClassifier\endcsname\trClassifier
\renewcommand{\trClassifier}{\BookTerm{T0110}{\csname BookOriginal@trClassifier\endcsname}}
\expandafter\def\csname BookIndex@T0111\endcsname{yubiaoshide@余表示的（corepresented）}
\expandafter\let\csname BookOriginal@trCorepresented\endcsname\trCorepresented
\renewcommand{\trCorepresented}{\BookTerm{T0111}{\csname BookOriginal@trCorepresented\endcsname}}
\expandafter\def\csname BookIndex@T0112\endcsname{danchunyingshe@单纯映射（simplicial map）}
\expandafter\let\csname BookOriginal@trSimplicialMap\endcsname\trSimplicialMap
\renewcommand{\trSimplicialMap}{\BookTerm{T0112}{\csname BookOriginal@trSimplicialMap\endcsname}}
\expandafter\def\csname BookIndex@T0113\endcsname{qunpei@∞-群胚（\ensuremath{\infty}-groupoid）}
\expandafter\let\csname BookOriginal@trInfinityGroupoid\endcsname\trInfinityGroupoid
\renewcommand{\trInfinityGroupoid}{\BookTerm{T0113}{\csname BookOriginal@trInfinityGroupoid\endcsname}}
\expandafter\def\csname BookIndex@T0114\endcsname{qiege@切割（cut）}
\expandafter\let\csname BookOriginal@trCut\endcsname\trCut
\renewcommand{\trCut}{\BookTerm{T0114}{\csname BookOriginal@trCut\endcsname}}
\expandafter\def\csname BookIndex@T0115\endcsname{qiepian@切片（slice）}
\expandafter\let\csname BookOriginal@trSlice\endcsname\trSlice
\renewcommand{\trSlice}{\BookTerm{T0115}{\csname BookOriginal@trSlice\endcsname}}
\expandafter\def\csname BookIndex@T0116\endcsname{zhui@锥（cone）}
\expandafter\let\csname BookOriginal@trCone\endcsname\trCone
\renewcommand{\trCone}{\BookTerm{T0116}{\csname BookOriginal@trCone\endcsname}}
\expandafter\def\csname BookIndex@T0117\endcsname{xianxingyouxuji@线性有序集（linearly ordered set）}
\expandafter\let\csname BookOriginal@trLinearlyOrderedSet\endcsname\trLinearlyOrderedSet
\renewcommand{\trLinearlyOrderedSet}{\BookTerm{T0117}{\csname BookOriginal@trLinearlyOrderedSet\endcsname}}
\expandafter\def\csname BookIndex@T0118\endcsname{yiwanghanzi@遗忘函子（forgetful functor）}
\expandafter\let\csname BookOriginal@trForgetfulFunctor\endcsname\trForgetfulFunctor
\renewcommand{\trForgetfulFunctor}{\BookTerm{T0118}{\csname BookOriginal@trForgetfulFunctor\endcsname}}
\expandafter\def\csname BookIndex@T0119\endcsname{fanbianhanzi@反变函子（contravariant functor）}
\expandafter\let\csname BookOriginal@trContravariantFunctor\endcsname\trContravariantFunctor
\renewcommand{\trContravariantFunctor}{\BookTerm{T0119}{\csname BookOriginal@trContravariantFunctor\endcsname}}
\expandafter\def\csname BookIndex@T0120\endcsname{kebiaoshide@可表示的（representable）}
\expandafter\let\csname BookOriginal@trRepresentable\endcsname\trRepresentable
\renewcommand{\trRepresentable}{\BookTerm{T0120}{\csname BookOriginal@trRepresentable\endcsname}}
\expandafter\def\csname BookIndex@T0121\endcsname{wanbeide@完备的（complete）}
\expandafter\let\csname BookOriginal@trComplete\endcsname\trComplete
\renewcommand{\trComplete}{\BookTerm{T0121}{\csname BookOriginal@trComplete\endcsname}}
\expandafter\def\csname BookIndex@T0122\endcsname{yuwanbeide@余完备的（cocomplete）}
\expandafter\let\csname BookOriginal@trCocomplete\endcsname\trCocomplete
\renewcommand{\trCocomplete}{\BookTerm{T0122}{\csname BookOriginal@trCocomplete\endcsname}}
\expandafter\def\csname BookIndex@T0123\endcsname{yudenghuazi@余等化子（coequalizer）}
\expandafter\let\csname BookOriginal@trCoequalizer\endcsname\trCoequalizer
\renewcommand{\trCoequalizer}{\BookTerm{T0123}{\csname BookOriginal@trCoequalizer\endcsname}}
\expandafter\def\csname BookIndex@T0124\endcsname{danwei@单位（unit）}
\expandafter\let\csname BookOriginal@trUnit\endcsname\trUnit
\renewcommand{\trUnit}{\BookTerm{T0124}{\csname BookOriginal@trUnit\endcsname}}
\expandafter\def\csname BookIndex@T0125\endcsname{yudanwei@余单位（counit）}
\expandafter\let\csname BookOriginal@trCounit\endcsname\trCounit
\renewcommand{\trCounit}{\BookTerm{T0125}{\csname BookOriginal@trCounit\endcsname}}
\expandafter\def\csname BookIndex@T0126\endcsname{liantong@连通（connected）}
\expandafter\let\csname BookOriginal@trConnected\endcsname\trConnected
\renewcommand{\trConnected}{\BookTerm{T0126}{\csname BookOriginal@trConnected\endcsname}}
\expandafter\def\csname BookIndex@T0127\endcsname{baohelei@饱和类（saturated class）}
\expandafter\let\csname BookOriginal@trSaturatedClass\endcsname\trSaturatedClass
\renewcommand{\trSaturatedClass}{\BookTerm{T0127}{\csname BookOriginal@trSaturatedClass\endcsname}}
\expandafter\def\csname BookIndex@T0128\endcsname{baoshoude@保守的（conservative）}
\expandafter\let\csname BookOriginal@trConservative\endcsname\trConservative
\renewcommand{\trConservative}{\BookTerm{T0128}{\csname BookOriginal@trConservative\endcsname}}
\expandafter\def\csname BookIndex@T0129\endcsname{xiangfanhanzi@相反函子（opposite functor）}
\expandafter\let\csname BookOriginal@trOppositeFunctor\endcsname\trOppositeFunctor
\renewcommand{\trOppositeFunctor}{\BookTerm{T0129}{\csname BookOriginal@trOppositeFunctor\endcsname}}
\expandafter\def\csname BookIndex@T0130\endcsname{fanchoudengjia@范畴等价（categorical equivalence）}
\expandafter\let\csname BookOriginal@trCategoricalEquivalence\endcsname\trCategoricalEquivalence
\renewcommand{\trCategoricalEquivalence}{\BookTerm{T0130}{\csname BookOriginal@trCategoricalEquivalence\endcsname}}
\expandafter\def\csname BookIndex@T0131\endcsname{zirandengjia@自然等价（natural equivalence）}
\expandafter\let\csname BookOriginal@trNaturalEquivalence\endcsname\trNaturalEquivalence
\renewcommand{\trNaturalEquivalence}{\BookTerm{T0131}{\csname BookOriginal@trNaturalEquivalence\endcsname}}
\expandafter\def\csname BookIndex@T0132\endcsname{dengjiaxianweihua@等价纤维化（isofibration）}
\expandafter\let\csname BookOriginal@trIsofibration\endcsname\trIsofibration
\renewcommand{\trIsofibration}{\BookTerm{T0132}{\csname BookOriginal@trIsofibration\endcsname}}
\expandafter\def\csname BookIndex@T0133\endcsname{zhudian@逐点（pointwise）}
\expandafter\let\csname BookOriginal@trPointwise\endcsname\trPointwise
\renewcommand{\trPointwise}{\BookTerm{T0133}{\csname BookOriginal@trPointwise\endcsname}}
\expandafter\def\csname BookIndex@T0134\endcsname{nijixian@逆极限（inverse limit）}
\expandafter\let\csname BookOriginal@trInverseLimit\endcsname\trInverseLimit
\renewcommand{\trInverseLimit}{\BookTerm{T0134}{\csname BookOriginal@trInverseLimit\endcsname}}
\expandafter\def\csname BookIndex@T0135\endcsname{xianwei@纤维（fibre）}
\expandafter\let\csname BookOriginal@trFibre\endcsname\trFibre
\renewcommand{\trFibre}{\BookTerm{T0135}{\csname BookOriginal@trFibre\endcsname}}
\expandafter\def\csname BookIndex@T0136\endcsname{tongchangfanchou@通常范畴（ordinary category）}
\expandafter\let\csname BookOriginal@trOrdinaryCategory\endcsname\trOrdinaryCategory
\renewcommand{\trOrdinaryCategory}{\BookTerm{T0136}{\csname BookOriginal@trOrdinaryCategory\endcsname}}
\expandafter\def\csname BookIndex@T0137\endcsname{sanquerxingzhi@三取二性质（3-for-2 property）}
\expandafter\let\csname BookOriginal@trThreeForTwoProperty\endcsname\trThreeForTwoProperty
\renewcommand{\trThreeForTwoProperty}{\BookTerm{T0137}{\csname BookOriginal@trThreeForTwoProperty\endcsname}}
\expandafter\def\csname BookIndex@T0138\endcsname{bilu@闭路（loop）}
\expandafter\let\csname BookOriginal@trLoop\endcsname\trLoop
\renewcommand{\trLoop}{\BookTerm{T0138}{\csname BookOriginal@trLoop\endcsname}}
\expandafter\def\csname BookIndex@T0139\endcsname{2baoqiang@2-胞腔（2-cell）}
\expandafter\let\csname BookOriginal@trTwoCell\endcsname\trTwoCell
\renewcommand{\trTwoCell}{\BookTerm{T0139}{\csname BookOriginal@trTwoCell\endcsname}}
\expandafter\def\csname BookIndex@T0140\endcsname{kenide@可逆的（invertible）}
\expandafter\let\csname BookOriginal@trInvertible\endcsname\trInvertible
\renewcommand{\trInvertible}{\BookTerm{T0140}{\csname BookOriginal@trInvertible\endcsname}}
\expandafter\def\csname BookIndex@T0141\endcsname{tonglunfanchou@同伦范畴（homotopy category）}
\expandafter\let\csname BookOriginal@trHomotopyCategory\endcsname\trHomotopyCategory
\renewcommand{\trHomotopyCategory}{\BookTerm{T0141}{\csname BookOriginal@trHomotopyCategory\endcsname}}
\expandafter\def\csname BookIndex@T0142\endcsname{xianweixingduixiang@纤维性对象（fibrant object）}
\expandafter\let\csname BookOriginal@trFibrantObject\endcsname\trFibrantObject
\renewcommand{\trFibrantObject}{\BookTerm{T0142}{\csname BookOriginal@trFibrantObject\endcsname}}
\expandafter\def\csname BookIndex@T0143\endcsname{xianweilie@纤维列（fibre sequence）}
\expandafter\let\csname BookOriginal@trFibreSequence\endcsname\trFibreSequence
\renewcommand{\trFibreSequence}{\BookTerm{T0143}{\csname BookOriginal@trFibreSequence\endcsname}}
\expandafter\def\csname BookIndex@T0144\endcsname{benzhixiang@本质像（essential image）}
\expandafter\let\csname BookOriginal@trEssentialImage\endcsname\trEssentialImage
\renewcommand{\trEssentialImage}{\BookTerm{T0144}{\csname BookOriginal@trEssentialImage\endcsname}}
\expandafter\def\csname BookIndex@T0145\endcsname{fanxingzhi@泛性质（universal property）}
\expandafter\let\csname BookOriginal@trUniversalProperty\endcsname\trUniversalProperty
\renewcommand{\trUniversalProperty}{\BookTerm{T0145}{\csname BookOriginal@trUniversalProperty\endcsname}}
\expandafter\def\csname BookIndex@T0146\endcsname{luxianweihua@路纤维化（path fibration）}
\expandafter\let\csname BookOriginal@trPathFibration\endcsname\trPathFibration
\renewcommand{\trPathFibration}{\BookTerm{T0146}{\csname BookOriginal@trPathFibration\endcsname}}
\expandafter\def\csname BookIndex@T0147\endcsname{luguo@滤过（filtration）}
\expandafter\let\csname BookOriginal@trFiltration\endcsname\trFiltration
\renewcommand{\trFiltration}{\BookTerm{T0147}{\csname BookOriginal@trFiltration\endcsname}}
\expandafter\def\csname BookIndex@T0148\endcsname{kesuode@可缩的（contractible）}
\expandafter\let\csname BookOriginal@trContractible\endcsname\trContractible
\renewcommand{\trContractible}{\BookTerm{T0148}{\csname BookOriginal@trContractible\endcsname}}
\expandafter\def\csname BookIndex@T0149\endcsname{geluotengdikegouzao@格罗滕迪克构造（Grothendieck construction）}
\expandafter\let\csname BookOriginal@trGrothendieckConstruction\endcsname\trGrothendieckConstruction
\renewcommand{\trGrothendieckConstruction}{\BookTerm{T0149}{\csname BookOriginal@trGrothendieckConstruction\endcsname}}
\expandafter\def\csname BookIndex@T0150\endcsname{daibiaojizuolingdiaode@带标记左零调的（marked left anodyne）}
\expandafter\let\csname BookOriginal@trMarkedLeftAnodyne\endcsname\trMarkedLeftAnodyne
\renewcommand{\trMarkedLeftAnodyne}{\BookTerm{T0150}{\csname BookOriginal@trMarkedLeftAnodyne\endcsname}}
\expandafter\def\csname BookIndex@T0151\endcsname{jubuyudikaerde@局部余笛卡尔的（locally cocartesian）}
\expandafter\let\csname BookOriginal@trLocallyCocartesian\endcsname\trLocallyCocartesian
\renewcommand{\trLocallyCocartesian}{\BookTerm{T0151}{\csname BookOriginal@trLocallyCocartesian\endcsname}}
\expandafter\def\csname BookIndex@T0152\endcsname{daibiaojidengjia@带标记等价（marked equivalence）}
\expandafter\let\csname BookOriginal@trMarkedEquivalence\endcsname\trMarkedEquivalence
\renewcommand{\trMarkedEquivalence}{\BookTerm{T0152}{\csname BookOriginal@trMarkedEquivalence\endcsname}}
\expandafter\def\csname BookIndex@T0153\endcsname{daibiaojibian@带标记边（marked edge）}
\expandafter\let\csname BookOriginal@trMarkedEdge\endcsname\trMarkedEdge
\renewcommand{\trMarkedEdge}{\BookTerm{T0153}{\csname BookOriginal@trMarkedEdge\endcsname}}
\expandafter\def\csname BookIndex@T0154\endcsname{dikaerbide@笛卡尔闭的（cartesian closed）}
\expandafter\let\csname BookOriginal@trCartesianClosed\endcsname\trCartesianClosed
\renewcommand{\trCartesianClosed}{\BookTerm{T0154}{\csname BookOriginal@trCartesianClosed\endcsname}}
\expandafter\def\csname BookIndex@T0155\endcsname{danchunmoxingfanchou@单纯模型范畴（simplicial model category）}
\expandafter\let\csname BookOriginal@trSimplicialModelCategory\endcsname\trSimplicialModelCategory
\renewcommand{\trSimplicialModelCategory}{\BookTerm{T0155}{\csname BookOriginal@trSimplicialModelCategory\endcsname}}
\expandafter\def\csname BookIndex@T0156\endcsname{shangxianweihua@上纤维化（cofibration）}
\expandafter\let\csname BookOriginal@trCofibration\endcsname\trCofibration
\renewcommand{\trCofibration}{\BookTerm{T0156}{\csname BookOriginal@trCofibration\endcsname}}
\expandafter\def\csname BookIndex@T0157\endcsname{shangxianweide@上纤维的（cofibrant）}
\expandafter\let\csname BookOriginal@trCofibrant\endcsname\trCofibrant
\renewcommand{\trCofibrant}{\BookTerm{T0157}{\csname BookOriginal@trCofibrant\endcsname}}
\expandafter\def\csname BookIndex@T0158\endcsname{shidingdianyingshe@始顶点映射（initial vertex map）}
\expandafter\let\csname BookOriginal@trInitialVertexMap\endcsname\trInitialVertexMap
\renewcommand{\trInitialVertexMap}{\BookTerm{T0158}{\csname BookOriginal@trInitialVertexMap\endcsname}}
\expandafter\def\csname BookIndex@T0159\endcsname{fanyudikaerxianweihua@泛余笛卡尔纤维化（universal cocartesian fibration）}
\expandafter\let\csname BookOriginal@trUniversalCocartesianFibration\endcsname\trUniversalCocartesianFibration
\renewcommand{\trUniversalCocartesianFibration}{\BookTerm{T0159}{\csname BookOriginal@trUniversalCocartesianFibration\endcsname}}
\expandafter\def\csname BookIndex@T0160\endcsname{tonglundikaerde@同伦笛卡尔的（homotopy cartesian）}
\expandafter\let\csname BookOriginal@trHomotopyCartesian\endcsname\trHomotopyCartesian
\renewcommand{\trHomotopyCartesian}{\BookTerm{T0160}{\csname BookOriginal@trHomotopyCartesian\endcsname}}
\expandafter\def\csname BookIndex@T0161\endcsname{changzhihanzi@常值函子（constant functor）}
\expandafter\let\csname BookOriginal@trConstantFunctor\endcsname\trConstantFunctor
\renewcommand{\trConstantFunctor}{\BookTerm{T0161}{\csname BookOriginal@trConstantFunctor\endcsname}}
\expandafter\def\csname BookIndex@T0162\endcsname{xiangfanfanchou@相反范畴（opposite category）}
\expandafter\let\csname BookOriginal@trOppositeCategory\endcsname\trOppositeCategory
\renewcommand{\trOppositeCategory}{\BookTerm{T0162}{\csname BookOriginal@trOppositeCategory\endcsname}}
\expandafter\def\csname BookIndex@T0163\endcsname{kuozhangwenti@扩张问题（extension problem）}
\expandafter\let\csname BookOriginal@trExtensionProblem\endcsname\trExtensionProblem
\renewcommand{\trExtensionProblem}{\BookTerm{T0163}{\csname BookOriginal@trExtensionProblem\endcsname}}
\expandafter\def\csname BookIndex@T0164\endcsname{shuangbiande@双变的（bivariant）}
\expandafter\let\csname BookOriginal@trBivariant\endcsname\trBivariant
\renewcommand{\trBivariant}{\BookTerm{T0164}{\csname BookOriginal@trBivariant\endcsname}}
\expandafter\def\csname BookIndex@T0165\endcsname{yuji@余积（coproduct）}
\expandafter\let\csname BookOriginal@trCoproduct\endcsname\trCoproduct
\renewcommand{\trCoproduct}{\BookTerm{T0165}{\csname BookOriginal@trCoproduct\endcsname}}
\expandafter\def\csname BookIndex@T0166\endcsname{ji@积（product）}
\expandafter\let\csname BookOriginal@trProduct\endcsname\trProduct
\renewcommand{\trProduct}{\BookTerm{T0166}{\csname BookOriginal@trProduct\endcsname}}
\expandafter\def\csname BookIndex@T0167\endcsname{lisanfanchou@离散范畴（discrete category）}
\expandafter\let\csname BookOriginal@trDiscreteCategory\endcsname\trDiscreteCategory
\renewcommand{\trDiscreteCategory}{\BookTerm{T0167}{\csname BookOriginal@trDiscreteCategory\endcsname}}
\expandafter\def\csname BookIndex@T0168\endcsname{biaoshiduixiang@表示对象（representing object）}
\expandafter\let\csname BookOriginal@trRepresentingObject\endcsname\trRepresentingObject
\renewcommand{\trRepresentingObject}{\BookTerm{T0168}{\csname BookOriginal@trRepresentingObject\endcsname}}
\expandafter\def\csname BookIndex@T0169\endcsname{ruokesuode@弱可缩的（weakly contractible）}
\expandafter\let\csname BookOriginal@trWeaklyContractible\endcsname\trWeaklyContractible
\renewcommand{\trWeaklyContractible}{\BookTerm{T0169}{\csname BookOriginal@trWeaklyContractible\endcsname}}
\expandafter\def\csname BookIndex@T0170\endcsname{gongweixing@共尾性（cofinality）}
\expandafter\let\csname BookOriginal@trCofinality\endcsname\trCofinality
\renewcommand{\trCofinality}{\BookTerm{T0170}{\csname BookOriginal@trCofinality\endcsname}}
\expandafter\def\csname BookIndex@T0171\endcsname{guanghuade@光滑的（smooth）}
\expandafter\let\csname BookOriginal@trSmooth\endcsname\trSmooth
\renewcommand{\trSmooth}{\BookTerm{T0171}{\csname BookOriginal@trSmooth\endcsname}}
\expandafter\def\csname BookIndex@T0172\endcsname{zhengchangde@正常的（proper）}
\expandafter\let\csname BookOriginal@trProper\endcsname\trProper
\renewcommand{\trProper}{\BookTerm{T0172}{\csname BookOriginal@trProper\endcsname}}
\expandafter\def\csname BookIndex@T0173\endcsname{fanguanghuade@泛光滑的（universally smooth）}
\expandafter\let\csname BookOriginal@trUniversallySmooth\endcsname\trUniversallySmooth
\renewcommand{\trUniversallySmooth}{\BookTerm{T0173}{\csname BookOriginal@trUniversallySmooth\endcsname}}
\expandafter\def\csname BookIndex@T0174\endcsname{fanzhengchangde@泛正常的（universally proper）}
\expandafter\let\csname BookOriginal@trUniversallyProper\endcsname\trUniversallyProper
\renewcommand{\trUniversallyProper}{\BookTerm{T0174}{\csname BookOriginal@trUniversallyProper\endcsname}}
\expandafter\def\csname BookIndex@T0175\endcsname{shuangdikaerxianweihua@双笛卡尔纤维化（bicartesian fibration）}
\expandafter\let\csname BookOriginal@trBicartesianFibration\endcsname\trBicartesianFibration
\renewcommand{\trBicartesianFibration}{\BookTerm{T0175}{\csname BookOriginal@trBicartesianFibration\endcsname}}
\expandafter\def\csname BookIndex@T0176\endcsname{shuangziranbianhuan@双自然变换（binatural transformation）}
\expandafter\let\csname BookOriginal@trBinaturalTransformation\endcsname\trBinaturalTransformation
\renewcommand{\trBinaturalTransformation}{\BookTerm{T0176}{\csname BookOriginal@trBinaturalTransformation\endcsname}}
\expandafter\def\csname BookIndex@T0177\endcsname{duijiaohanzi@对角函子（diagonal functor）}
\expandafter\let\csname BookOriginal@trDiagonalFunctor\endcsname\trDiagonalFunctor
\renewcommand{\trDiagonalFunctor}{\BookTerm{T0177}{\csname BookOriginal@trDiagonalFunctor\endcsname}}
\expandafter\def\csname BookIndex@T0178\endcsname{benzhixiaode@本质小的（essentially small）}
\expandafter\let\csname BookOriginal@trEssentiallySmall\endcsname\trEssentiallySmall
\renewcommand{\trEssentiallySmall}{\BookTerm{T0178}{\csname BookOriginal@trEssentiallySmall\endcsname}}
\expandafter\def\csname BookIndex@T0179\endcsname{jubuxiaode@局部小的（locally small）}
\expandafter\let\csname BookOriginal@trLocallySmall\endcsname\trLocallySmall
\renewcommand{\trLocallySmall}{\BookTerm{T0179}{\csname BookOriginal@trLocallySmall\endcsname}}
\expandafter\def\csname BookIndex@T0180\endcsname{ruoshide@弱始的（weakly initial）}
\expandafter\let\csname BookOriginal@trWeaklyInitial\endcsname\trWeaklyInitial
\renewcommand{\trWeaklyInitial}{\BookTerm{T0180}{\csname BookOriginal@trWeaklyInitial\endcsname}}
\expandafter\def\csname BookIndex@T0181\endcsname{feituihuade@非退化的（non degenerate）}
\expandafter\let\csname BookOriginal@trNonDegenerate\endcsname\trNonDegenerate
\renewcommand{\trNonDegenerate}{\BookTerm{T0181}{\csname BookOriginal@trNonDegenerate\endcsname}}
\expandafter\def\csname BookIndex@T0182\endcsname{shousuoyingshe@收缩映射（retraction）}
\expandafter\let\csname BookOriginal@trRetraction\endcsname\trRetraction
\renewcommand{\trRetraction}{\BookTerm{T0182}{\csname BookOriginal@trRetraction\endcsname}}
\expandafter\def\csname BookIndex@T0183\endcsname{daidiande@带点的（pointed）}
\expandafter\let\csname BookOriginal@trPointed\endcsname\trPointed
\renewcommand{\trPointed}{\BookTerm{T0183}{\csname BookOriginal@trPointed\endcsname}}
\expandafter\def\csname BookIndex@T0184\endcsname{shizhui@始锥（initial cone）}
\expandafter\let\csname BookOriginal@trInitialCone\endcsname\trInitialCone
\renewcommand{\trInitialCone}{\BookTerm{T0184}{\csname BookOriginal@trInitialCone\endcsname}}
\expandafter\def\csname BookIndex@T0185\endcsname{kebiaoxiande@可表现的（presentable）}
\expandafter\let\csname BookOriginal@trPresentable\endcsname\trPresentable
\renewcommand{\trPresentable}{\BookTerm{T0185}{\csname BookOriginal@trPresentable\endcsname}}
\expandafter\def\csname BookIndex@T0186\endcsname{kejide@可及的（accessible）}
\expandafter\let\csname BookOriginal@trAccessible\endcsname\trAccessible
\renewcommand{\trAccessible}{\BookTerm{T0186}{\csname BookOriginal@trAccessible\endcsname}}
\expandafter\def\csname BookIndex@T0187\endcsname{ruozhongde@弱终的（weakly terminal）}
\expandafter\let\csname BookOriginal@trWeaklyTerminal\endcsname\trWeaklyTerminal
\renewcommand{\trWeaklyTerminal}{\BookTerm{T0187}{\csname BookOriginal@trWeaklyTerminal\endcsname}}
\expandafter\def\csname BookIndex@T0188\endcsname{yuan@源（source）}
\expandafter\let\csname BookOriginal@trSource\endcsname\trSource
\renewcommand{\trSource}{\BookTerm{T0188}{\csname BookOriginal@trSource\endcsname}}
\expandafter\def\csname BookIndex@T0189\endcsname{ba@靶（target）}
\expandafter\let\csname BookOriginal@trTarget\endcsname\trTarget
\renewcommand{\trTarget}{\BookTerm{T0189}{\csname BookOriginal@trTarget\endcsname}}
\expandafter\def\csname BookIndex@T0190\endcsname{pu@谱（spectrum）}
\expandafter\let\csname BookOriginal@trSpectrum\endcsname\trSpectrum
\renewcommand{\trSpectrum}{\BookTerm{T0190}{\csname BookOriginal@trSpectrum\endcsname}}
\expandafter\def\csname BookIndex@T0191\endcsname{ceng@层（sheaf）}
\expandafter\let\csname BookOriginal@trSheaf\endcsname\trSheaf
\renewcommand{\trSheaf}{\BookTerm{T0191}{\csname BookOriginal@trSheaf\endcsname}}
\expandafter\def\csname BookIndex@T0192\endcsname{xiajiang@下降（descent）}
\expandafter\let\csname BookOriginal@trDescent\endcsname\trDescent
\renewcommand{\trDescent}{\BookTerm{T0192}{\csname BookOriginal@trDescent\endcsname}}
\expandafter\def\csname BookIndex@T0193\endcsname{qunwanbeihua@群完备化（group completion）}
\expandafter\let\csname BookOriginal@trGroupCompletion\endcsname\trGroupCompletion
\renewcommand{\trGroupCompletion}{\BookTerm{T0193}{\csname BookOriginal@trGroupCompletion\endcsname}}
\expandafter\def\csname BookIndex@T0194\endcsname{gaixing@概形（scheme）}
\expandafter\let\csname BookOriginal@trScheme\endcsname\trScheme
\renewcommand{\trScheme}{\BookTerm{T0194}{\csname BookOriginal@trScheme\endcsname}}
\expandafter\def\csname BookIndex@T0195\endcsname{daochufanchou@导出范畴（derived category）}
\expandafter\let\csname BookOriginal@trDerivedCategory\endcsname\trDerivedCategory
\renewcommand{\trDerivedCategory}{\BookTerm{T0195}{\csname BookOriginal@trDerivedCategory\endcsname}}
\expandafter\def\csname BookIndex@T0196\endcsname{daji@打积（smash product）}
\expandafter\let\csname BookOriginal@trSmashProduct\endcsname\trSmashProduct
\renewcommand{\trSmashProduct}{\BookTerm{T0196}{\csname BookOriginal@trSmashProduct\endcsname}}
\expandafter\def\csname BookIndex@T0197\endcsname{mutide@母题的（motivic）}
\expandafter\let\csname BookOriginal@trMotivic\endcsname\trMotivic
\renewcommand{\trMotivic}{\BookTerm{T0197}{\csname BookOriginal@trMotivic\endcsname}}
\expandafter\def\csname BookIndex@T0198\endcsname{jubuxitong@局部系统（local system）}
\expandafter\let\csname BookOriginal@trLocalSystem\endcsname\trLocalSystem
\renewcommand{\trLocalSystem}{\BookTerm{T0198}{\csname BookOriginal@trLocalSystem\endcsname}}
\expandafter\def\csname BookIndex@T0199\endcsname{gebulilun@割补理论（surgery theory）}
\expandafter\let\csname BookOriginal@trSurgeryTheory\endcsname\trSurgeryTheory
\renewcommand{\trSurgeryTheory}{\BookTerm{T0199}{\csname BookOriginal@trSurgeryTheory\endcsname}}
\expandafter\def\csname BookIndex@T0200\endcsname{manzifanchou@满子范畴（full subcategory）}
\expandafter\let\csname BookOriginal@trFullSubcategory\endcsname\trFullSubcategory
\renewcommand{\trFullSubcategory}{\BookTerm{T0200}{\csname BookOriginal@trFullSubcategory\endcsname}}
\expandafter\def\csname BookIndex@T0201\endcsname{fenleiyingshe@分类映射（classifying map）}
\expandafter\let\csname BookOriginal@trClassifyingMap\endcsname\trClassifyingMap
\renewcommand{\trClassifyingMap}{\BookTerm{T0201}{\csname BookOriginal@trClassifyingMap\endcsname}}
\expandafter\def\csname BookIndex@T0202\endcsname{tonglunjiashe@同伦假设（homotopy hypothesis）}
\expandafter\let\csname BookOriginal@trHomotopyHypothesis\endcsname\trHomotopyHypothesis
\renewcommand{\trHomotopyHypothesis}{\BookTerm{T0202}{\csname BookOriginal@trHomotopyHypothesis\endcsname}}
\expandafter\def\csname BookIndex@T0203\endcsname{ziyouzhulitonggou@自由驻立同构（free-standing isomorphism）}
\expandafter\let\csname BookOriginal@trFreeStandingIsomorphism\endcsname\trFreeStandingIsomorphism
\renewcommand{\trFreeStandingIsomorphism}{\BookTerm{T0203}{\csname BookOriginal@trFreeStandingIsomorphism\endcsname}}
\expandafter\def\csname BookIndex@T0204\endcsname{jubuchangzhide@局部常值的（locally constant）}
\expandafter\let\csname BookOriginal@trLocallyConstant\endcsname\trLocallyConstant
\renewcommand{\trLocallyConstant}{\BookTerm{T0204}{\csname BookOriginal@trLocallyConstant\endcsname}}
\expandafter\def\csname BookIndex@T0205\endcsname{aiermite@埃尔米特（hermitian）}
\expandafter\let\csname BookOriginal@trHermitian\endcsname\trHermitian
\renewcommand{\trHermitian}{\BookTerm{T0205}{\csname BookOriginal@trHermitian\endcsname}}
\expandafter\def\csname BookIndex@T0206\endcsname{kehechengde@可合成的（composable）}
\expandafter\let\csname BookOriginal@trComposable\endcsname\trComposable
\renewcommand{\trComposable}{\BookTerm{T0206}{\csname BookOriginal@trComposable\endcsname}}
\expandafter\def\csname BookIndex@T0207\endcsname{fanchou@∞-范畴（\ensuremath{\infty} categories）}
\expandafter\let\csname BookOriginal@trInfinityCategories\endcsname\trInfinityCategories
\renewcommand{\trInfinityCategories}{\BookTerm{T0207}{\csname BookOriginal@trInfinityCategories\endcsname}}
\expandafter\def\csname BookIndex@T0208\endcsname{fanchou@范畴（categories）}
\expandafter\let\csname BookOriginal@trCategories\endcsname\trCategories
\renewcommand{\trCategories}{\BookTerm{T0208}{\csname BookOriginal@trCategories\endcsname}}
\expandafter\def\csname BookIndex@T0209\endcsname{danchunji@单纯集（simplicial sets）}
\expandafter\let\csname BookOriginal@trSimplicialSets\endcsname\trSimplicialSets
\renewcommand{\trSimplicialSets}{\BookTerm{T0209}{\csname BookOriginal@trSimplicialSets\endcsname}}
\expandafter\def\csname BookIndex@T0210\endcsname{danchunfanchou@单纯范畴（simplicial categories）}
\expandafter\let\csname BookOriginal@trSimplicialCategories\endcsname\trSimplicialCategories
\renewcommand{\trSimplicialCategories}{\BookTerm{T0210}{\csname BookOriginal@trSimplicialCategories\endcsname}}
\expandafter\def\csname BookIndex@T0211\endcsname{hechengqi@合成器（weak composer）}
\expandafter\let\csname BookOriginal@trWeakComposer\endcsname\trWeakComposer
\renewcommand{\trWeakComposer}{\BookTerm{T0211}{\csname BookOriginal@trWeakComposer\endcsname}}
\expandafter\def\csname BookIndex@T0212\endcsname{baohebibao@饱和闭包（saturation/saturated closure）}
\expandafter\def\csname BookIndex@T0213\endcsname{bianjie@边界（boundary）}
\expandafter\def\csname BookIndex@T0214\endcsname{bousfield@Bousfield（Bousfield）}
\expandafter\def\csname BookIndex@T0215\endcsname{danchunhengdengshi@单纯恒等式（simplicial identities）}
\expandafter\def\csname BookIndex@T0216\endcsname{danchunshenjing@单纯神经（simplicial nerve）}
\expandafter\def\csname BookIndex@T0217\endcsname{2fanchou@2-范畴（2-category）}
\expandafter\def\csname BookIndex@T0218\endcsname{waijiao@外角（outer horn）}
\expandafter\def\csname BookIndex@T0219\endcsname{weihanzi@伪函子（pseudo-functor）}
\expandafter\def\csname BookIndex@T0220\endcsname{yaobanhanzi@幺半函子（monoidal functor）}
\expandafter\def\csname BookIndex@T0221\endcsname{youtishengxingzhi@右提升性质（RLP）}
\expandafter\def\csname BookIndex@T0222\endcsname{zuotishengxingzhi@左提升性质（LLP）}
\expandafter\def\csname BookIndex@T0223\endcsname{dehuaier@德怀尔（William Dwyer）}
\expandafter\def\csname BookIndex@T0224\endcsname{huaitehai@怀特海（J. H. C. Whitehead）}
\expandafter\def\csname BookIndex@T0225\endcsname{shidiwensen@史蒂文森（Daniel Stevenson）}
\expandafter\def\csname BookIndex@T0226\endcsname{pangjialai@庞加莱（Henri Poincaré）}
\expandafter\def\csname BookIndex@T0227\endcsname{adiya@阿蒂亚（Michael Atiyah）}
\makeatother

